\section{Related Work}

\noindent\textbf{Urban perception and design cues.}
Crowdsourced pairwise judgement datasets have enabled vision models for
perceived safety, wealth, and liveliness at
scale~\cite{salesses2013,naik2014,dubey2016,quintana2025}, but these
models are not designed to answer mechanistic scene-level questions.
Urban-design theory links perceived safety to surveillance, active fontages,
upkeep, and legibility~\cite{jacobs1961,newman1972,loewen1993}, and
empirical street-view studies confirm roles for greenery, street-facing
windows, maintenance, and lighting~\cite{li2015greenery,denadai2016lively,jiang2018brokenwindows,portnov2020lighting}. However, none provides scene-level tests of localised, validity-audited
single-lever interventions.

\noindent\textbf{Interpretability and counterfactual explanations.}
Saliency and SHAP-style methods are known to be visually plausible yet
unfaithful~\cite{adebayo2018}; concept-based methods such as
TCAV~\cite{kim2018} remain correlational unless paired with explicit
interventions. Counterfactual visual explanations~\cite{goyal2019} and
causal evaluation benchmarks~\cite{melistas2024} provide stronger tools
but typically target model logits, focusing on the effect of a single feature on the model output rather than the holistic perceptual experience of the scene which can be affected by multiple single visual changes.

\noindent\textbf{Urban counterfactuals and diffusion-based editing.}
Law et al.~\cite{law2023} demonstrate plausible counterfactuals for
urban image regressors, establishing that generative edits can serve as
explanations in urban analytics but does not explicitly identify which edits are plausible for a given scene. UrbanPhysicalDisorder-4K~\cite{vera2025} brings counterfactual reasoning
into urban safety via annotated disorder features, but operates in
feature space rather than through localised image edits with explicit
validity auditing. Zhao et al.~\cite{zhao2026} propose
perception-guided street-view generation, but target scene optimisation
rather than scene-level interventional counterfactual analysis.
Diffusion-based editors such as SDEdit~\cite{meng2022},
InstructPix2Pix~\cite{brooks2023}, and
Prompt-to-Prompt~\cite{hertz2023prompt} make localised edits
increasingly feasible, but can introduce non-target drift that demands
explicit validity auditing.

\noindent\textbf{Positioning.}
The diffusion-editing literature above optimises for editing
fidelity, how faithfully a model executes an instruction, whereas our work sits at this intersection: rather than attributing a scene-level judgement to pixels or concepts alone, we search over structured, editor-feasible single-lever edits and retain only those that pass explicit validity auditing.